% One round of a toy 16-bit SPN: key addition, 4 parallel 4-bit S-boxes,
% bit permutation P(i) = 4i mod 15 (i < 15), P(15) = 15 (PRESENT-style).
\documentclass[border=3pt]{standalone}
\usepackage[T1]{fontenc}
\usepackage{lmodern}
\usepackage{amsmath}
\usepackage{crypto-tikz}
\begin{document}
\begin{tikzpicture}[x=4.2mm, y=-1cm]
  % input bits
  \foreach \i in {0,...,15} {
    \draw[wire] (\i,0) -- (\i,0.55);
    \node[xorop] (k\i) at (\i,0.7) {};
    \draw[wire] (\i,0.85) -- (\i,1.25);
  }
  \node[op, anchor=east] at (-0.8,0.7) {$\oplus\, K_r$};
  \node[op, anchor=east] at (-0.8,0)  {$x^{(r)}$};
  % S-box layer
  \foreach \s in {0,...,3}
    \node[sbox, minimum width=15mm] at ({4*\s+1.5},1.55) {$S$};
  \node[op, anchor=east] at (-0.8,1.55) {S-layer};
  % permutation layer
  \foreach \i in {0,...,15} {
    \pgfmathtruncatemacro{\j}{\i==15 ? 15 : mod(4*\i,15)}
    \draw[wire] (\i,1.85) -- (\i,2.05) -- (\j,3.05) -- (\j,3.3);
  }
  \node[op, anchor=east] at (-0.8,2.55) {$P$};
  \node[op, anchor=east] at (-0.8,3.3) {$x^{(r+1)}$};
  \draw[brace] (-0.2,-0.25) -- node[op, above=5pt] {16-bit state} (15.2,-0.25);
\end{tikzpicture}
\end{document}
